\section{Introdução e objetivos}
\label{sec:objetivos}

Em um amplificador operacional com realimentação negativa, os ganhos
de malha fechada são definidos por razões de resistores, o que permite
somar e subtrair sinais com coeficientes
arbitrários. O roteiro solicita o projeto de uma rede
desse tipo, com coeficientes definidos pelos dígitos da matrícula.
Para a matrícula 21101175 tem-se $A_7{=}2$, $A_6{=}1$, $A_5{=}1$,
$A_4{=}0$, $A_3{=}1$, $A_2{=}1$, $A_1{=}7$ e $A_0{=}5$. Substituindo
os dígitos, as funções que cada bloco deve implementar são:
%
\[ v_{out1} = (5-2A_1)v_{in1} + \frac{A_2}{5} - 3v_{in2} + \frac{6-A_3}{2}v_{in3}
            = -9v_{in1} + 0{,}2 - 3v_{in2} + 2{,}5v_{in3} \]
\[ v_{out2} = (A_4+2)v_{in4} + A_5[4v_{in5} - (10-A_6)v_{in6}]
            = 2v_{in4} + 4v_{in5} - 9v_{in6} \]
\[ v_{out} = \frac{G}{2}(v_{out2}-v_{out1}) = 12{,}5\,(v_{out2}-v_{out1}),
   \qquad G = 10A_7+A_0 = 25 \]
%
O termo $A_2/5 = 0{,}2$~V do bloco A é uma constante, ou seja, um
nível DC somado ao sinal. A composição dos três blocos resulta na
função global (deduzida na Seção~\ref{sec:composicao})
%
\[ v_{out} = 112{,}5v_{in1} + 37{,}5v_{in2} - 31{,}25v_{in3}
           + 25v_{in4} + 50v_{in5} - 112{,}5v_{in6} - 2{,}5. \]

O amplificador operacional adotado é o LM741 em modelo ideal nível 1,
com alimentação simétrica de $\pm$\SI{15}{\volt}. O roteiro prevê o
uso do Micro-Cap; neste trabalho as simulações foram realizadas no
LTspice. A troca de simulador não altera o procedimento: o modelo
nível 1 foi reproduzido no LTspice com os mesmos parâmetros
(Seção~\ref{sec:modelo}) e a alimentação simétrica, que o Micro-Cap
insere automaticamente, foi incluída de forma explícita.

Os objetivos do experimento são: (i) projetar o bloco A utilizando
somente somadores e subtratores, o bloco B com topologia livre e o
bloco C com um amplificador de instrumentação, descrevendo e
justificando as escolhas de projeto; (ii) validar por simulação o
efeito individual de cada entrada sobre a saída e a implementação da
função global; e (iii) comparar o resultado do modelo ideal com o
macromodelo comercial do LM741.
